% Sample Applications, Chapter 6 -- Invoice Aging (annex Sec 8.3).
% Values replayed through the shipped application (examples/aging6.i23):
% @DATE serial numbers (os/i123/core/eval.c), the Date 2 display
% (os/i123/core/fmt.c), a difference of two dates as days, @IF.
% Figure screens/sa6.k123.
% GATES: test-i123-formula test-pubs-i123-examples test-pubs-screens
\Gates{test-i123-formula, test-pubs-i123-examples, test-pubs-screens}
\chapter{Invoice Aging}

The accounts receivable clerk reviews unpaid invoices at the end of each
month and calls the customers whose invoices are more than sixty days old.
This worksheet works out how many days each invoice has been outstanding and
marks the late ones.

\section{What You Enter}
\begin{itemize}
\item The date of the review
\item The date of each invoice
\end{itemize}

\section{Entering the Data}
\ix{invoice aging}\ix{dates}\ix{"@IF}
Dates are serial numbers (Tutorial, Lesson~22).  Type each date as its serial
number; \texttt{@DATE(92,6,30)}, for example, gives 33785 for the thirtieth of
June.  Type \type{As of} in A1 and the date of the review, \type{33785}, in
B1.  Type the headings \type{"Date}, \type{"Days} and \type{"Status} in
A2..C2 and the invoice dates \type{33704}, \type{33746} and \type{33767} in
A3..A5.  In B3 enter \type{+\$b\$1-a3}, and in C3 enter
\type{@if(b3>60,"Late","OK")}.  Copy B3..C3 to B4 and to B5.  Finally, format
A3..A5 and B1 with \slashcmd{/Range Format Date 2 (DD-MMM)}, which shows the day and
month.
\FlairShot[fig:sa6]{sa6}{Days outstanding and the late invoices}

\begin{LotTable}{p{30mm}p{24mm}p{70mm}}
\LotHead{\textbf{Item} & \textbf{Address} & \textbf{Meaning}}
As of & B1 & The date of the review, as a serial number\\
Date & A3..A5 & The date of each invoice, as a serial number\\
Days & B3..B5 & The review date minus the invoice date: the days outstanding\\
Status & C3..C5 & \texttt{Late} when the invoice is more than sixty days
  old, otherwise \texttt{OK}\LotEnd
\end{LotTable}

\section{Interpreting the Results}
\LotExample{examples/aging6}
The first four lines show that the serial numbers you typed are the ones
\fn{@DATE} gives.  The formula in B3 subtracts one date from another, so it
yields a number of days.  The reference \texttt{\$B\$1} is absolute, so every
copy subtracts from the same review date, while \texttt{A3} moves down with
the row.  \fn{@IF} compares the days with sixty and returns a label.

\section{Changing the Template}
To review a different day, type its serial number in B1.  To mark invoices
more than thirty days old as well, replace the formula in C3 with
\begin{center}\type{@IF(B3>60,"Late",@IF(B3>30,"Due","OK"))}\end{center}
and copy it to C4 and C5.
